\chapter{Database.java}
\section{Background}

\paragraph{}
The Marks Administration System would work with .csv files that would be stored in a database.
However, working with actual databases, would result in using external libraries
like JDBC. To work around that, a folder was used in place of a porper database.
The folder stores the marks in a marks.csv file, and the methods in Database.java 
all work together to perform the basic functions of a databse, that being Create, Update, Read and Delete.
The file also needed to have functions that prevented the program from crashing due to commas in fields or strings, 
split lines and rows into fields and collumns, along with functions that convert values from doubles to strings, and in vice versa.

\section{Supporter Functions}

\paragraph{}
There are 5 smaller functions that are used throughout four larger functions in Database.java.
The functions are:
\begin{verbatim}
private String escapeCSV(String value)
private String[] splitCSV(String line)
private String markToString(Double mark)
private Double stringToMark(String value)
\end{verbatim}

\paragraph{}
The escapeCSV method prevents the CSV from being broken by commas, by wrapping a field in 
quotes if it contains a comma. The splitCSV method splits the lines and rows of the CSV
into fields and collumns, with respect to commas. A CSV row is merely text split by ommas, so 
commas inside a fiel would be identical to a collumn seperator. Without escapeCSV, if a lecturer adds a name and surname seperated by a comma inside of a field,
then the saveData function, will put both name and surname in more than 5 fields. 
Without splitCSV, the data in each field is not seperated properly which can return errors if a name 
is in a field where a double should be.

\paragraph{}
In the process of gathering data and populating files with it, mark values needs to be 
converted from double to string and in vice versa, depending on the current step in the method.
If data is being taken in from a CSV file, every double is a string and must be converted into a double 
so that it can be accepted as a parameter in program methods.

\section{saveData function}
\paragraph{}
The saveData function performs the Create and Update feature of the database, whilst simulatenously doing deletion.
The algorithm of the method creates a list of lines that are populated with a header and the data 
retrieved from the studentManager object that serves as its only parameter. The studentmanager will have Information 
regarding students, their names, ids and marks. All this information, and retreived in a for loop where 
new student objects are made, and for every student, their information is given via the manager.
All this data is used to populate the rows with student information.

\paragraph{}
At the bottom of the method is a a try-catch block, made incase of failure. If the file to be saved to according a 
Database object is not found, then the directory will be made and the file will be written to 
with the information stored in the list of lines. This unintentionally results in the saveData method also 
deleting/overwritting existing data in a file, provided the data wasn't loaded previously onto the studentManager using 
loadData

\section{loadData function}
\paragraph{}
loadData is a public StudentManager method that uses assessment component parameters.
The reason for this is because this method is what retains the data on files that are to be worked on 
by lecturers, and returns a studentManager object. This object contains a List that reads all the lines in the file 
located at the data folder, and in a for loop, it retrieves student information, loads it into a student object and 
adds the student to student manager. An if statement above prevents an infinity loop, checking if the lines are blank and 
ceasing operations if they are. The for loop also begins in the row below the header.

\section{importCSV function}
\paragraph{}
One of the CSV managing files, importCSV, uses the assessment components as parameters. 
Initially, it used the components and a path, but the path was removed. If the 
function is to be used in the main file, the path cannot be set outside the function. 
It needs to prompt the user to enter the path of the source file to be imported. 
The source file is then read line by line, until it reaches an empty line.

\paragraph{}
An issue with importing files is that external files may not follow the order of the program.
Collumns may be named different names or be in different places. If this isn't accounted for then the program risks breaking.
The function has a block of code made for dealing with this. The code 
fills an array with the words in the header, and creates a String-Integer Map 
to track collumn indexes. For each name found, it is trimmed and reduced to lowercase, then put into the map with its collumn index. 
The header strings are also checked against a requiredheader array, which stores every expected header,
and throws an exception if there isn't an expected header.
The for loop through the lines does what loadData does. It loads all information into 
student objects and adds to a manager, then returns it. 

\section{exportCSV function}
\paragraph{}
This function is a near exact copy of saveData. This is because to export\
needs the data to be filled on a file that must be created at the destination directory.
This is what saveData does but locally for the program in the data folder.
Unlike saveData, loadData doesn't need to be called before an exported file is made.
The function is designed to create a parent directory either way, and the scanner inside the function, 
allows for the user to name the exported file anything they want.

\paragraph{}
A problem that may arise is whether or there is an already existing file of the same name 
at the destination. A check might need to be implemented in future, after the prototype stage.
 